\documentclass[11pt,a4paper]{article}
\usepackage[margin=2.5cm]{geometry}
\usepackage{amsmath,amssymb}
\usepackage{booktabs}
\usepackage{listings}
\usepackage{xcolor}
\usepackage{hyperref}
\usepackage{enumitem}

\lstset{
  basicstyle=\ttfamily\footnotesize,
  breaklines=true,
  columns=flexible,
  showstringspaces=false,
}

\title{SLAC Light--Dark POMDP: Debugging Report}
\author{}
\date{\today}

\begin{document}
\maketitle

\section{Starting point}

The codebase contains two main training scripts that implement SLAC (Stochastic
Latent Actor--Critic) with C51 distributional Q-learning for a partially
observable light--dark navigation task:

\begin{itemize}[leftmargin=*]
  \item \texttt{SLAC\_light\_dark\_POMDP\_C51\_MI\_old\_working.py} --- the
        reference implementation that converges: world-model loss reaches
        $\sim 10^{-4}$ during the main loop, and the agent achieves
        \texttt{eval/success\_reach\_rate} up to $0.38$ at 500k env steps.
  \item \texttt{SLAC\_light\_dark\_POMDP\_C51\_MI\_new.py} --- the active
        development version. At the start of this debugging session it did not
        learn: \texttt{success\_strict\_rate} stayed at $0$ throughout
        training, eval return remained near $-1.9$.
\end{itemize}

Both scripts share the same world-model code
(\texttt{SLAC\_Agent\_deterministic\_tabular.py}), the same C51 agent
(\texttt{SLAC\_Agent\_D3QN\_tabular.py}), and the same environment
(\texttt{envs/Light\_dark\_POMDP\_flags.py}). The architecture is a two-level
latent hierarchy: \(z_1 \in \mathbb{R}^{32}\) (stochastic, fast-changing) and
\(z_2 \in \mathbb{R}^{256}\) (effectively deterministic, slow-changing).
Training proceeds in three phases:
\begin{enumerate}[leftmargin=*]
  \item Bootstrap (10k random-action steps).
  \item World-model pretraining (100k steps): first half MSE-only on the
        decoder reconstruction; second half adds KL, latent transfer MSE, and
        pixel transfer MSE.
  \item Joint world-model and C51 learning (500k env steps) with an
        information-theoretic intrinsic reward.
\end{enumerate}

\section{Diagnostic findings}

\subsection{Frozen \texorpdfstring{$z_1$}{z1} standard deviation}

Posterior and prior networks for $z_1$ both initialize with
$\sigma \approx 1.87$ (from the softplus parameterization). The KL term
$\mathrm{KL}(q\,\|\,p)$ has gradient that cancels when both distributions have
identical scale, so the standard deviation does not move. Because the decoder
can route information through $z_2$ alone (which is high-dimensional and
quasi-deterministic), the reconstruction MSE term has no incentive to push
$z_1.\sigma$ down either. The result: $z_1$ carries no information about
the observation.

\subsection{World-model loss does not recover after the joint phase begins}

Pretraining MSE drops to $\sim 6 \times 10^{-6}$ in the MSE-only first half,
then jumps abruptly to $\sim 0.19$ when KL, latent transfer MSE, and pixel
transfer MSE turn on simultaneously, and plateaus at $\sim 0.14$ for the
remainder of pretraining. In the main loop the world-model loss stays at
$\sim 0.35$. In contrast, the old working file reaches
\texttt{train/world\_model\_loss} $\approx 1{-}4 \times 10^{-4}$ in the main
loop --- roughly three orders of magnitude lower.

\subsection{Latent collapse and flat Q-values}

In the new file:
\begin{itemize}[leftmargin=*]
  \item Posterior $z_2$ std collapses to $\sim 5 \times 10^{-3}$ --- $z_2$ is
        effectively deterministic.
  \item Per-action Q-value spread \texttt{diag/q\_taken\_std\_mean} stays at
        $\sim 5 \times 10^{-4}$. The C51 critic cannot differentiate between
        the 25 discrete actions.
  \item Mutual-information bonus \texttt{intrinsic/raw\_mean} is
        $\sim 10^{-5}$. With $z_2$ near-deterministic and the C51 distribution
        invariant across posterior samples, $I(Z;R\mid a) \approx 0$ as
        expected.
\end{itemize}

\section{Changes attempted}

\subsection{NaN crashes in the inference path}

Early runs crashed inside \texttt{latent2\_prior} during the Bayes-filter
rollout with all-\texttt{NaN} \texttt{loc}. Diagnostics showed weights were
finite at the moment of failure --- the NaN arose from intermediate-activation
overflow during a forward pass. Two clamps were added to
\texttt{MLPCompoundNormalSimo.forward} in
\texttt{SLAC\_Agent\_deterministic\_tabular.py}:
\begin{lstlisting}[language=Python]
x  = x.clamp(min=-30.0, max=30.0)                # input
loc = out[..., :latent_size].clamp(min=-20, max=20)
scale_diag = F.softplus(...).clamp(max=2.0) + 1e-5
\end{lstlisting}
Combined with a gradient-norm clip on the world-model optimizer, NaN crashes
stopped. The clip was later raised from $1.0$ to $100$ once measured
gradient norms during the joint phase ($\sim 0.5$--$1.0$) showed it was
binding and likely throttling joint-phase recovery; however, even at clip
$=100$ the world-model loss did not improve.

\subsection{Decoder input}

For some intermediate runs the decoder was forced to reconstruct from $z_1$
alone (\texttt{z\_2} replaced by zeros), in an attempt to break the
$z_1$-frozen-std symmetry. This unfroze $z_1$
($\sigma_{z_1} \approx 0.83$ instead of $1.87$) but introduced inference-time
drift in $z_2$, causing the NaN crashes mentioned above. The decoder was
reverted to receive the real $z_2$, matching the old working file.

\subsection{Intrinsic-reward mode}

The new file introduced an
\texttt{intrinsic\_mode} switch with two new variants:
\begin{itemize}[leftmargin=*]
  \item \texttt{posterior\_reduction}: a posterior-entropy-reduction bonus,
        unbounded; observed \texttt{raw\_mean} $\sim 70$--$90$, which the
        running auto-scaler crushed to near-zero.
  \item \texttt{action\_disagreement}: posterior-sample argmax disagreement.
  \item \texttt{mi\_bonus} (restored): the true mutual information
        $I(Z;R\mid a) = H(\bar{p}(R\mid z,a)) - \mathbb{E}_z\,H(p(R\mid z,a))$
        used in the old working file. Bounded by $\log 51 \approx 3.9$ nats.
\end{itemize}
The default was switched to \texttt{mi\_bonus} since the old working file uses
this formulation. With the world-model still broken, however, the MI bonus is
$\sim 10^{-5}$ regardless --- not a usable signal.

\subsection{Hyperparameters tested}

\begin{itemize}[leftmargin=*]
  \item \texttt{sigma\_dark}: $3.0 \to 2.0$ (matching the old working file's
        env construction).
  \item \texttt{latent2\_size}: briefly $256 \to 32$ and reverted to $256$.
  \item \texttt{world\_model\_update\_frequency}: $32 \to 16$ (matching the
        old working file's main-loop update cadence).
\end{itemize}

None of these individually closed the gap.

\section{Where we stand}

The two scripts are nominally near-identical for the pretraining phase:
\begin{itemize}[leftmargin=*]
  \item Same imports, same model file, same optimizer, same learning rate
        ($10^{-4}$), batch size ($128$), sequence length ($8$), latent sizes
        ($32$ + $256$), pretrain budget ($100{,}000$ steps), seed ($42$).
  \item Same environment configuration ($\sigma_{\mathrm{dark}}=2.0$,
        $\sigma_{\mathrm{light}}=0.05$, \texttt{noisy\_goal\_obs=False}).
  \item Same \texttt{compute\_loss} structure: MSE + adaptive-weighted KL with
        free bits + $0.1\,\cdot$ \texttt{latent\_tf\_mse} +
        \texttt{pix\_tf\_mse}.
\end{itemize}
The only remaining textual differences in \texttt{compute\_loss} are
cosmetic, plus a minor \texttt{adaptive\_w} floor change from $0$ to $0.05$
that does not bind in the observed regime.

Despite this, pretraining outcomes differ by orders of magnitude:

\begin{center}
\begin{tabular}{lll}
\toprule
phase & old working file & new file \\
\midrule
end of MSE-only pretrain  & --- (not logged)             & $\sim 6 \times 10^{-6}$ \\
end of joint pretrain     & low (inferred)               & $\sim 0.14$             \\
main-loop world model     & $\sim 1\text{--}4 \times 10^{-4}$ & $\sim 0.35$        \\
\texttt{success\_reach\_rate} @ 500k & $0.32$                       & $0.04$ \\
\bottomrule
\end{tabular}
\end{center}

The new file's main script has $\sim 400$ lines of cumulative uncommitted diff
from its initial commit, including added diagnostics, the intrinsic-mode
switch, and assorted refactors. Static inspection has not located a single
change that explains the world-model divergence. The remaining hypothesis is
a subtle drift in the main loop body or in a helper that we have not yet
isolated.

\section{Proposed next step}

Rather than continuing the change-by-change bisection, replace
\texttt{SLAC\_light\_dark\_POMDP\_C51\_MI\_new.py} with a copy of
\texttt{SLAC\_light\_dark\_POMDP\_C51\_MI\_old\_working.py} and verify the
baseline reproduces. Then conservatively re-introduce, one group at a time:
\begin{enumerate}[leftmargin=*]
  \item the \texttt{intrinsic\_mode} switch and the \texttt{mi\_bonus} branch
        using the pre-existing \texttt{compute\_mi\_bonus},
  \item pretrain wandb logging,
  \item the \texttt{diag/*} diagnostic computations.
\end{enumerate}
After each re-introduction, run end-to-end and confirm
\texttt{train/world\_model\_loss} reaches $\sim 10^{-4}$ in the main loop
before adding the next group.

\end{document}
